\paragraph{Direct TDX Attestation Bottleneck.}
In a conventional DCAP-based deployment, each end-user attestation request
must carry a fresh nonce, and that nonce must be embedded into a newly
generated TDX quote. Consequently, $N$ independent users induce $N$ hardware
quote-generation operations. Although PCCS/DCAP removes the Intel Trust
Authority network dependency, it does not remove the per-request quote cost:
the TDX module must produce a nonce-bound report and the quoting path must
serialize access to the attestation key material, making quote generation the
dominant scalability bottleneck while quote verification remains comparatively
cheap. As user demand increases, the TDX VM therefore behaves as a queued
single-service point whose throughput is bounded by hardware quote generation.
Vordr addresses this limitation by moving TDX attestation off the critical
path: a WEN attests the TDX VM once per refresh interval and then serves many
end users from cached verified state, amortizing one expensive TDX attestation
across a large number of end-user attestation responses.

